% Part: lambda-calculus
% Chapter: introduction
% Section: free-variables

\documentclass[../../../include/open-logic-section]{subfiles}

\begin{document}

\olfileid{lam}{syn}{fv}
\olsection{آزاد متغیر}

لامبڈا کیلکولس فنکشوناں بارے اے، تے فنکشن لامبڈا تجرید نال
بندے نیں۔ بدیہی خیال ایہ اے کہ $\lambd[x][M]$ اوہ فنکشن اے
جیہدیاں قدراں~$M$ توں ملدیاں نیں، جدوں فنکشن دی دلیل
نوں~$x$ دی قدر بنایا جاوے۔ پر~$x$ دی~$M$ وچ ہر آمد متعلقہ
نہیں۔ جے $M$ وچ اک ہور تجرید $\lambd[x][N]$ اے، تاں $x$
دیاں~$N$ وچلیاں آمدواں $\lambd[x][N]$ نال متعلق نیں،
$\lambd[x][M]$ نال نہیں۔ سو لامبڈا تجرید $\lambd[x]$،
$\lambd[x][M]$ دے اندر، $x$ دیاں~$M$ وچ اوہناں آمدواں
نوں \emph{پابند کردی اے} جیہڑیاں پہلاں توں کسے ہور
لامبڈا تجرید نال پابند نہیں، یعنی~$x$ دیاں~$M$ وچ
\emph{آزاد} آمدواں۔

\begin{defn}[دائرۂ اثر]
جے $\lambd[x][M]$ ٹرم~$N$ دے اندر آوے، تاں~$N$ دی متعلقہ
آمد~$\lambd[x]$ دا \emph{دائرۂ اثر} اے۔
\end{defn}

\begin{defn}[آزاد تے پابند آمد]
متغیر $x$ دی ٹرم~$M$ وچ آمد \emph{آزاد} اے جے اوہ کسے
$\lambd[x]$ دے دائرۂ اثر وچ نہیں، تے ہور صورت وچ \emph{پابند} اے۔
متغیر~$x$ دی $\lambd[x][M]$ وچ آمد مُڈھلے $\lambd[x]$ نال
اُدوں تے صرف اُدوں پابند اے جدوں $x$ دی اوہ آمد~$M$ وچ آزاد اے۔
\end{defn}

\begin{ex}
$\lambd[x][x y]$ وچ $x$ تے $y$ دوویں $\lambd[x]$ دے دائرۂ
اثر وچ نیں، سو $x$، $\lambd[x]$ نال پابند اے۔ کیوں جے $y$
کسے $\lambd[y]$ دے دائرۂ اثر وچ نہیں، اوہ آزاد اے۔
$\lambd[x][x x]$ وچ $x$ دیاں دوویں آمدواں~$\lambd[x]$ نال پابند
نیں، کیوں جے دوویں $xx$ وچ آزاد نیں۔ $((\lambd[x][xx])x)$ وچ~$x$
دی آخری آمد آزاد اے، کیوں جے اوہ کسے $\lambd[x]$ دے دائرۂ اثر وچ
نہیں۔ $\lambd[x][(\lambd[x][x])x]$ وچ پہلے $\lambd[x]$ دا دائرۂ
اثر $(\lambd[x][x])x$ اے، تے دوجے $\lambd[x]$ دا دائرۂ اثر~$x$
دی آخری توں پہلی آمد اے۔ $(\lambd[x][x])x$ وچ $x$ دی آخری آمد آزاد
اے، تے آخری توں پہلی پابند اے۔ ایس لئی $x$ دی آخری توں پہلی آمد،
$\lambd[x][(\lambd[x][x])x]$ وچ دوجے $\lambd[x]$ نال پابند اے،
تے آخری آمد پہلے $\lambd[x]$ نال پابند اے۔
\end{ex}

ٹرم $P$ وچ متغیراں دیاں ساریاں آمدواں پرکھ کے اسیں آزاد
متغیراں دا سیٹ لبھ سکدے آں۔ ایہ سیٹ $\FV{P}$ نال لکھدے آں، تے
اوہدی قدرتی تعریف ایہ اے:

\begin{defn}[ٹرم دے آزاد متغیر] \ollabel{def:fv}
ٹرم دے \emph{آزاد متغیراں} دا سیٹ استقرا نال ایویں متعین کردے آں:
\begin{enumerate}
  \item $\FV{x} = \{x\}$ \ollabel{def:fv1}
  \item $\FV{\lambd[x][N]} = \FV{N} \setminus \{x\}$ \ollabel{def:fv2}
  \item $\FV{PQ} = \FV{P} \cup \FV{Q}$ \ollabel{def:fv3}
\end{enumerate}
\end{defn}

\begin{prob}
\begin{enumerate}
  \item $\lambd[g]$ تے دوواں $\lambd[x]$ دے دائرۂ اثر ایس ٹرم وچ لبھو:
    $\lambd[g][(\lambd[x][g (x x)]) \lambd[x][g (x x)]]$۔
  \item $\lambd[g][(\lambd[x][g (x x)]) \lambd[x][g (x x)]]$ وچ
    کی متغیراں دیاں ساریاں آمدواں پابند نیں؟ ہر اک نوں کیہڑی تجرید
    پابند کردی اے؟
  \item $\FV{\lambd[x][(\lambd[y][(\lambd[z][x y]) z]) y]}$ دسو۔
\end{enumerate}
\end{prob}

\begin{explain}
آزاد متغیر باہرلی دنیا، یعنی \emph{ماحول}، ول اک حوالے وانگ اے۔
جیہڑے ٹرم وچ آزاد متغیر ہون، اوہنوں ادھورا متعین ٹرم سمجھ سکدے آں،
کیوں جے اوہدا ورتاؤ ماحول نوں کیویں مقرر کردے آں، ایس اُتے منحصر اے۔
مثال لئی، ٹرم $\lambd[x][f x]$ اک دلیل~$x$ لیندا اے تے اوس دلیل
اُتے $f$ دا نتیجہ واپس کردا اے۔ ایتھے متغیر~$f$ آزاد اے۔ ٹرم دی
ایہ قدر اوہدے ماحول اُتے، خاص کر اوس ماحول وچ $f$ دی قدر اُتے،
منحصر اے۔

جے ایس ٹرم اُتے تجرید لائیے، تاں $\lambd[f][\lambd[x][f
    x]]$ ملدا اے۔ ایہ ٹرم ہُن ماحول دے متغیر $f$ اُتے منحصر نہیں۔
ہُن ایہ اوہ فنکشن ظاہر کردا اے جیہڑا دو دلیل لیندا اے تے پہلی
نوں دوجی اُتے لا کے نتیجہ واپس کردا اے۔ ماحول وچ $f$ نوں بدلن
نال ایس ٹرم دے ورتاؤ اُتے کوئی اثر نہیں پئے گا، کیوں جے ٹرم
صرف اوہو قدر ورتے گا جیہڑی دلیل دے طور اُتے دِتّی جاوے گی۔
ماحول وچ $f$ دی قدر نہیں ورتے گا۔
\end{explain}

\begin{defn}[بند ٹرم، کمبینیٹر]
جیہڑے ٹرم وچ کوئی آزاد متغیر نہ ہووے، اوہنوں \emph{بند ٹرم}
یا \emph{کمبینیٹر} آکھدے آں۔
\end{defn}

\begin{lem}\ollabel{lem:fv}
\begin{enumerate}
  \item \ollabel{lem:fv-abs} جے $y \neq x$، تاں $y \in
    \FV{\lambd[x][N]}$ اُدوں تے صرف اُدوں جدوں $y \in \FV{N}$۔
  \item \ollabel{lem:fv-app} $y \in \FV{PQ}$ اُدوں تے صرف اُدوں
    جدوں $y \in \FV{P}$ یا $y \in \FV{Q}$۔
\end{enumerate}
\end{lem}

\begin{proof}
مشق۔
\end{proof}

\begin{prob}
\olref[lam][syn][fv]{lem:fv} ثابت کرو۔
\end{prob}

\end{document}
